problem:  
Let \((M^n, g)\) be a compact, simply connected Riemannian manifold of dimension \(n\) with nonnegative sectional curvature, diameter at most \(D\), and satisfying \(\| \operatorname{Rm}_g \|_{C^0} \leq 1\). Assume that there exists a smooth function \(f : M \to \mathbb{R}\) and a constant \(\lambda \in \mathbb{R}\) such that  
\[
\operatorname{Ric}_g + \nabla^2 f = \lambda g + E, \qquad \|E\|_{C^2} \leq \varepsilon.
\]  
Does there exist, for each dimension \(n\) and diameter bound \(D\), a constant \(\varepsilon_0(n, D) > 0\) such that if \(\varepsilon < \varepsilon_0\), then \((M,g)\) is \(C^{2,\alpha}\)-close (modulo diffeomorphism and scaling) to an Einstein metric \((M, g_E)\)? Equivalently, does quantitative almost Einstein rigidity hold within the class of manifolds with bounded geometry?  

Why is it a "good" problem:  
This formulation introduces precise geometric control—bounded curvature, diameter, and regularity of the error tensor—so the question becomes logically sound and anchored in the framework of geometric compactness and partial differential equations. It asks whether small analytic deviation from the Einstein condition implies geometric and topological rigidity, a delicate and largely unresolved stability question in Riemannian geometry.  

The problem bridges analytic and geometric viewpoints: it studies the nonlinear elliptic PDE \(\operatorname{Ric}_g = \lambda g - \nabla^2 f + E\) through the lens of global geometric analysis and moduli-space stability theory. Its resolution would unify Ricci soliton stability, Gromov–Hausdorff almost-rigidity, and Einstein moduli space compactness into a quantitative and curvature-controlled framework.  

Moreover, the question isolates a profound conceptual frontier: whether “Einstein structure” is *stable under almost Einstein perturbations* in the \(C^2\) topology with bounded geometry. A positive answer would establish a dimension- and geometry-dependent rigidity gap generalizing classical sphere and hyperbolic rigidity phenomena, whereas a counterexample would reveal new moduli instability mechanisms within Ricci solitons.  

Thus, this problem is simple to state yet deeply intertwined with core ideas of geometric analysis, PDE stability, and global topology—embodying the essence of a "good" research-level problem.